%============================================================
% CHƯƠNG 5: ĐẠO LÀM NGƯỜI
%============================================================
\chapter{Đạo Làm Người (Being a Good Human)}

\begin{center}
  \textit{\color{truyenblue}``The true measure of any society can be found in how it treats its most vulnerable members.''}\\
  \textit{(Thước đo thực sự của một xã hội nằm ở cách nó đối xử với những thành viên dễ bị tổn thương nhất.)}\\[4pt]
  \small\color{truyengold}--- Mahatma Gandhi ---
\end{center}
\ngancach

\section{Đồng Tiền Lăn}
\begin{truyen}{Đồng Tiền Lăn}{Trung Thực}
\chuhoa{M}{ }ột người ăn xin đang chết đói nhặt được đồng tiền vàng rơi từ túi người lái xe sang.\\
\textit{(A starving beggar picked up a gold coin that had fallen from a wealthy driver's pocket.)}

Anh có thể mua đủ thứ ăn cho cả tháng. Nhưng anh \textbf{chạy theo xe hơi để trả lại}.\\
\textit{(He could buy food for a whole month. But he \textbf{ran after the car to return it}.)}

Người lái xe ngạc nhiên hỏi: ``Tại sao anh không giữ lại?''\\
\textit{(The driver asked in surprise: ``Why did you not keep it?'')}

``\textbf{Đói có thể giết thân xác tôi. Nhưng gian lận sẽ giết linh hồn tôi. Tôi chỉ mất một cái thôi}.''\\
\textit{(``\textbf{Hunger can kill my body. But dishonesty will kill my soul. I would rather lose only one of them}.'')}
\end{truyen}
\begin{baihoc}
  Lòng tự trọng không phải là xa xỉ phẩm của kẻ giàu. Nó là thứ phẩm giá thiêng liêng mà người nghèo nhất cũng có quyền và có sức gìn giữ.
\end{baihoc}

\section{Người Đếm Sao}
\begin{truyen}{Người Đếm Sao}{Lạc Quan}
\chuhoa{M}{ }ột người đàn ông vui vẻ ngồi giữa vũng bùn lầy và ngắm nhìn bầu trời đầy sao.\\
\textit{(A cheerful man sat right in the middle of a muddy puddle and gazed up at a starry sky.)}

Người đi qua chê cười: ``Ông điên à? Đang ngồi trong bùn mà cười được!''\\
\textit{(Passers-by sneered: ``Are you mad? Sitting in the mud and still smiling!'')}

Ông nhìn lên không ngừng mỉm cười: ``\textbf{Nếu tôi nhìn xuống, tôi chỉ thấy bùn. Tôi chọn nhìn lên bầu trời sao. Hai người cùng ngồi trong bùn, nhưng chỉ một người thấy các ngôi sao}.''\\
\textit{(He kept looking up and smiling: ``\textbf{If I look down, I only see mud. I choose to look up at the stars. Two people sit in the same mud, but only one sees the stars}.'')}
\end{truyen}
\begin{baihoc}
  Hoàn cảnh giống nhau có thể tạo ra trải nghiệm hoàn toàn khác nhau tùy thuộc vào góc nhìn của chúng ta. Lạc quan không phải là phủ nhận khó khăn mà là chọn nhìn xa hơn nó.
\end{baihoc}

\section{Người Trồng Cây}
\begin{truyen}{Người Trồng Cây}{Vô Tư}
\chuhoa{M}{ }ột bà lão già yếu dành cả đời trồng cây ăn quả dọc con đường làng.\\
\textit{(A frail old woman spent her whole life planting fruit trees along the village road.)}

Người hàng xóm cáu kỉnh: ``Bà già rồi, không ăn được đâu. Trồng làm chi vất vả.''\\
\textit{(A grumpy neighbour said: ``You are old, you will never get to eat them. Why bother?'')}

Bà khẽ đặt hạt mầm cuối cùng xuống đất: ``\textbf{Đúng vậy. Nhưng những người lữ hành mệt mỏi ngày mai sẽ được nghỉ mát dưới bóng cây này}.''\\
\textit{(The old woman gently placed the last seed in the ground: ``\textbf{Exactly. But the weary travellers of tomorrow will rest under the shade of this tree}.'')}

Khi bà qua đời, dọc con đường ấy \textbf{mát rượi và trĩu quả suốt bốn mùa}.\\
\textit{(When she passed away, that road was \textbf{cool and laden with fruit throughout all four seasons}.)}
\end{truyen}
\begin{baihoc}
  Đạo làm người đẹp nhất là sống không chỉ cho mình. Hành động gieo mầm cho thế hệ sau mà không cần gặt hái là biểu hiện cao nhất của lòng nhân ái.
\end{baihoc}

\section{Tiếng Chuông Gió}
\begin{truyen}{Tiếng Chuông Gió}{Tha Thứ}
\chuhoa{H}{ }ai nhà hàng xóm không nhìn mặt nhau suốt mười năm vì tranh chấp đất.\\
\textit{(Two neighbours had not looked at each other for ten years over a land dispute.)}

Một đêm bão lớn thổi vỡ chuông gió vườn nhà người kia, mảnh thủy tinh rơi vãi khắp sân.\\
\textit{(One stormy night, the wind shattered the other's garden wind chime, scattering glass shards across the yard.)}

Sáng hôm sau, người hàng xóm đầu tiên thức dậy sớm, lặng lẽ \textbf{quét sạch mảnh vỡ vì sợ trẻ con nhà kia bị đứt chân}.\\
\textit{(The next morning, the first neighbour woke early and quietly \textbf{swept up every shard, fearing the other's children might get cut}.)}

Người kia nhìn thấy từ cửa sổ, \textbf{ra mở cổng và gõ cửa nhà hàng xóm lần đầu tiên sau mười năm}.\\
\textit{(The other saw it from the window, \textbf{opened the gate and knocked on the neighbour's door for the first time in ten years}.)}
\end{truyen}
\begin{baihoc}
  Một hành động nhân ái nhỏ đôi khi có sức mạnh phá tan bức tường hận thù mà ngàn lời nói không làm được. Tính người thể hiện rõ nhất khi ta quan tâm đến người chưa kịp làm lành với mình.
\end{baihoc}

\section{Kẻ Không Tên}
\begin{truyen}{Kẻ Không Tên}{Lòng Tốt Vô Danh}
\chuhoa{M}{ }ỗi sáng thứ Hai, người lao công trường học luôn tìm thấy một ly cà phê nóng đặt trước cửa phòng làm việc.\\
\textit{(Every Monday morning, the school janitor always found a hot cup of coffee placed outside his workroom door.)}

Không ai biết của ai. Không có ghi chú.\\
\textit{(No one knew who left it. There was no note.)}

Sau ba năm, ông hỏi thăm nhiều người và mãi không tìm ra.\\
\textit{(After three years he asked around and never found out.)}

Nhưng mỗi sáng thứ Hai, \textbf{ông đi làm sớm hơn nửa tiếng và quét sạch hơn mọi khi} – không phải vì quy định, mà vì ai đó đã nhớ đến ông.\\
\textit{(But every Monday, \textbf{he arrived half an hour earlier and swept more thoroughly than usual} – not because of rules, but because someone had remembered him.)}
\end{truyen}
\begin{baihoc}
  Lòng tốt không cần danh tiếng mới có giá trị. Đôi khi một cử chỉ vô danh nhỏ bé tạo ra làn sóng nhân ái lớn gấp nhiều lần nguồn gốc ban đầu.
\end{baihoc}

\section{Người Đẩy Xe}
\begin{truyen}{Người Đẩy Xe}{Giúp Đỡ}
\chuhoa{T}{ }rời mưa nặng hạt, một người đàn ông mặc vest đắt tiền dừng lại giúp bà cụ đẩy xe hàng bị thụt bùn.\\
\textit{(In heavy rain, a man in an expensive suit stopped to help an old woman push her cart stuck in the mud.)}

Bộ vest ướt hết, bùn bắn khắp người. Người đi đường nhìn ngơ ngác.\\
\textit{(His suit was soaked, mud splattered everywhere. Passers-by stared in disbelief.)}

Khi xe ra khỏi bùn, bà cụ cảm ơn rối rít.\\
\textit{(When the cart was free from the mud, the old woman thanked him profusely.)}

Ông cười và đi tiếp, không bận tâm bộ vest.\\
\textit{(He smiled and walked on, unconcerned about the suit.)}

Tối về, ông kể chuyện đó cho con nghe như \textbf{một ngày bình thường. Đó là điều làm nó trở nên phi thường}.\\
\textit{(That evening he told his children about it as \textbf{an ordinary day. That is precisely what made it extraordinary}.)}
\end{truyen}
\begin{baihoc}
  Lòng tốt thực sự là thứ được làm như chuyện hiển nhiên, không phải như thành tích cần khoe. Người làm điều tốt như chuyện bình thường mới là người có nhân cách đã ngấm vào xương tủy.
\end{baihoc}

\section{Chỗ Ngồi Trên Tàu}
\begin{truyen}{Chỗ Ngồi Trên Tàu}{Nhường Nhịn}
\chuhoa{T}{ }rên chuyến tàu đêm chật kín người, một thanh niên đang ngồi đọc sách.\\
\textit{(On a packed night train, a young man was sitting and reading.)}

Một bà mẹ bồng con nhỏ đứng gần đó. Hành khách xung quanh đều giả vờ ngủ.\\
\textit{(A mother with a small child in her arms stood nearby. The surrounding passengers all pretended to be asleep.)}

Chàng trai gấp sách lại, đứng dậy nhường chỗ.\\
\textit{(The young man closed his book, stood up and offered his seat.)}

Bà mẹ cảm ơn, rồi thêm: ``Cảm ơn vì không giả vờ ngủ.''\\
\textit{(The mother thanked him, then added: ``Thank you for not pretending to be asleep.'')}

Chàng trai đứng suốt ba tiếng còn lại nhưng \textbf{đọc được thêm vài chục trang và ngủ ngon hơn đêm đó}.\\
\textit{(The young man stood for the remaining three hours but \textbf{read a few more dozen pages and slept better that night}.)}
\end{truyen}
\begin{baihoc}
  Làm một việc đúng khi người xung quanh đang quay mặt đi là hành động đòi hỏi dũng cảm thực sự. Lương tâm bình an là phần thưởng không thể mua được.
\end{baihoc}

\section{Lọ Muối}
\begin{truyen}{Lọ Muối}{Góc Nhìn}
\chuhoa{S}{ }ư phụ hỏi đệ tử: ``Con hãy hòa một thìa muối vào bát nước này rồi uống thử.''\\
\textit{(The master asked the student: ``Dissolve a spoonful of salt in this bowl of water and taste it.'')}

Đệ tử nhăn mặt: ``Mặn lắm thưa thầy.''\\
\textit{(The student grimaced: ``It is very salty, master.'')}

``Bây giờ thả thìa muối đó vào hồ nước trước sân và uống thử.''\\
\textit{(``Now drop that same spoonful of salt into the pond outside and taste it.'')}

``Ngọt lành thưa thầy.''\\
\textit{(``It tastes fresh, master.'')}

``\textbf{Muối khổ đau của cuộc đời không thay đổi. Chỉ có lòng ta rộng hay hẹp mới quyết định vị của nó}.''\\
\textit{(``\textbf{The salt of life's suffering does not change. Only whether our heart is wide or narrow determines its taste}.'')}
\end{truyen}
\begin{baihoc}
  Không thể thay đổi tất cả những gì xảy đến với mình, nhưng hoàn toàn có thể thay đổi dung lượng trái tim đón nhận nó. Rộng lòng là liều thuốc chữa lành mọi đau khổ.
\end{baihoc}

\section{Người Bán Hoa}
\begin{truyen}{Người Bán Hoa}{Tử Tế}
\chuhoa{M}{ }ột buổi sáng mưa, cô bé bán hoa đứng mãi không ai mua. Hoa sắp tàn.\\
\textit{(One rainy morning, the flower girl stood for a long time with no buyers. The flowers were starting to wilt.)}

Một người đàn ông đi qua, mua tất cả số hoa còn lại mặc dù không có ai để tặng.\\
\textit{(A man passed by and bought all the remaining flowers even though he had no one to give them to.)}

Cô bé hỏi: ``Chú mua hoa tặng ai vậy?''\\
\textit{(The girl asked: ``Who are you buying flowers for, sir?'')}

Ông nhìn những bông hoa sắp tàn và mỉm cười: ``\textbf{Tặng cho buổi sáng hôm nay. Để buổi sáng này không trở thành một buổi sáng u ám cho em}.''\\
\textit{(He looked at the wilting flowers and smiled: ``\textbf{For this morning. So this morning does not become a gloomy one for you}.'')}
\end{truyen}
\begin{baihoc}
  Sự tử tế không nhất thiết phải có mục đích cụ thể. Đôi khi chỉ cần giữ cho một buổi sáng của người khác không trở nên buồn bã là đã đủ ý nghĩa để hành động.
\end{baihoc}

\section{Gương Không Phản Chiếu}
\begin{truyen}{Gương Không Phản Chiếu}{Tự Nhìn Lại}
\chuhoa{M}{ }ột người hay chỉ trích mọi người xung quanh than thở với sư phụ: ``Sao con người thời nay tệ quá vậy thầy?''\\
\textit{(A person who often criticised everyone around complained to the master: ``Why are people so terrible these days?'')}

Sư phụ đưa cho anh ta chiếc gương: ``Nhìn vào đây và cho ta biết con thấy gì?''\\
\textit{(The master handed him a mirror: ``Look into this and tell me what you see.'')}

``Con thấy con thưa thầy.''\\
\textit{(``I see myself, master.'')}

``\textbf{Hãy nhớ lấy điều đó. Mỗi lần con chuẩn bị phán xét người khác, hãy nhìn vào gương trước}. Người ta chỉ thấy trong người khác thứ mà họ đã quen nhìn thấy trong chính mình.''\\
\textit{(``\textbf{Remember that. Every time you are about to judge someone, look in the mirror first}. People only see in others what they are already familiar with seeing in themselves.'')}
\end{truyen}
\begin{baihoc}
  Đạo làm người bắt đầu từ tự nhận thức. Trước khi phán xét thế giới, ta cần đủ dũng cảm để nhìn thẳng vào gương và đối diện với bản thân mình.
\end{baihoc}
